\section{Threat Model}
\label{sec:threat}

\paragraph{Assets.} OTPs; account credentials; phone numbers, email addresses, and postal addresses; payment and account identifiers; and any secret whose value is usable if disclosed.

\paragraph{Trusted.} The \texttt{system\_server} and the \SH{} service; the OS source pipelines (SMS/notification OTP classifier, Credential Manager, clipboard, wallet); and hardware-backed key storage (StrongBox, or a software keystore where hardware is unavailable).

\paragraph{Untrusted.} The agent and its model, which may be prompt-injected; third-party applications, one of which may render handle-like strings or send a fake OTP message; cloud model endpoints; and web content reached through an in-app WebView.

\paragraph{Attacker goals.} (G1) read a secret the agent handles; (G2) redirect a secret to an attacker-controlled destination; (G3) replay or reuse a secret; (G4) forge a handle.

\paragraph{Security goals.} For a task that legitimately requires a secret, \SH{} aims to guarantee: (S1) the raw value never enters the agent's observation or a non-bound destination; (S2) handles are unforgeable OS capabilities; (S3) resolved values are redacted from subsequent agent observation; (S4) single-use secrets cannot be reused. The utility goal is that login, 2FA, and checkout flows still complete.

\paragraph{Agent principals.} Mobile agents reach a device in more than one way, so the agent is not a single principal. We model three: (i) an accessibility/media-projection service of an application (application UID); (ii) an ADB/UiAutomation process, which many open runners use and which runs as the \texttt{shell} UID; and (iii) a remote agent that only receives serialized observations. The resolver authenticates the caller's UID at entry, and the redactor keys redaction to the principal at the screenshot and accessibility boundaries. If the agent has root, or the device is not SELinux-enforcing under an unrooted ADB, the redaction boundary is void and results are not meaningful; enforcement is therefore evaluated under \texttt{adb unroot} with SELinux enforcing.

\paragraph{Out of scope.} Semantic misuse \emph{inside} an allowed destination (e.g., an agent tricked into initiating a legitimate-but-unwanted action), WebView-origin attestation (fields in a WebView are rejected fail-closed rather than bound), and side channels outside the modeled observation boundary. These are residual risks and are discussed in the evaluation plan and limitations, not claimed to be eliminated.
